\documentclass[12pt, a4paper]{article}

\usepackage{amsmath, amssymb, amsthm}
\usepackage{geometry}
\usepackage{hyperref}
\usepackage{titlesec}
\usepackage{parskip}
\usepackage{microtype}
\usepackage{lmodern}
\usepackage{xcolor}
\usepackage{booktabs}
\usepackage{enumitem}

\geometry{margin=1.2in}

\definecolor{deepblue}{RGB}{20, 40, 100}
\definecolor{muted}{RGB}{90, 90, 90}

\hypersetup{
  colorlinks=true,
  linkcolor=deepblue,
  urlcolor=deepblue,
  citecolor=deepblue
}

\titleformat{\section}{\large\bfseries\color{deepblue}}{{\thesection}}{1em}{}
\titleformat{\subsection}{\normalsize\bfseries\color{deepblue}}{{\thesubsection}}{1em}{}

\theoremstyle{definition}
\newtheorem{definition}{Definition}
\newtheorem{question}{Research Question}

% Notation shortcuts
\newcommand{\R}{\mathbb{R}}
\newcommand{\N}{\mathbb{N}}
\newcommand{\q}{\mathbf{q}}
\newcommand{\Action}{\mathcal{A}}
\newcommand{\Lag}{\mathcal{L}}
\newcommand{\Sym}{\mathrm{Sym}}
\newcommand{\SO}{\mathrm{SO}}
\newcommand{\Z}{\mathbb{Z}}

\begin{document}

% ---------- TITLE ----------
\begin{center}
  {\LARGE \textbf{Three-Dimensional Choreographic Orbits:\\[4pt]
  A Point-Group Classification Programme}}\\[1.6em]
  {\large \textcolor{muted}{Research Design Document — v0.2}}\\[0.6em]
  {\normalsize \textcolor{muted}{e5bed97b3349 \quad\textbullet\quad March 2026}}
\end{center}

\vspace{1.2em}
\noindent\rule{\textwidth}{0.4pt}
\vspace{0.4em}

% ---------- ABSTRACT ----------
\begin{abstract}
\noindent
We propose a systematic programme to discover, classify, and characterise
three-dimensional equal-mass choreographic solutions to the Newtonian
$N$-body problem. A \emph{choreographic orbit} is a periodic solution in
which all $N$ bodies traverse a single closed curve, equally spaced in
phase. While the two-dimensional case has been surveyed extensively — most
notably by Sim\'o (2002) — fully three-dimensional choreographies remain
largely unexplored. Our programme organises the search space by the
three-dimensional point-group symmetry of the orbit, imports the
classification machinery of molecular symmetry into gravitational dynamics,
and characterises each discovered orbit by the knot type of its spacetime
trace. We expect the resulting catalogue to reveal a rich correspondence
between symmetry class, knot invariant, and action value, and to draw
explicit connections to the emerging literature on mechanically synthesised
molecular knots.
\end{abstract}

\noindent\rule{\textwidth}{0.4pt}
\vspace{1em}

% ---------- 1. BACKGROUND ----------
\section{Background and Motivation}

\subsection{The choreographic orbit}

Consider $N$ bodies of equal mass $m$ moving in $\R^3$ under mutual
Newtonian gravity. A \emph{choreographic solution} is a periodic orbit
$\q_i(t) \in \R^3$ ($i = 1, \ldots, N$, period $T$) such that there exists
a single closed curve $\gamma \subset \R^3$ and a phase shift $\tau = T/N$
with
\[
  \q_i(t) = \q_1\!\left(t + (i-1)\,\tau\right), \quad i = 1, \ldots, N.
\]
All bodies follow the same spatial curve; they are simply offset in time.
The solution is therefore entirely specified by the single curve
$\gamma$ and its parameterisation.

\subsection{Historical context}

The existence of non-trivial choreographic solutions was not established
rigorously until 2000, when Chenciner and Montgomery \cite{CM2000} proved
the existence of the figure-eight solution for $N = 3$ by showing it
minimises the Lagrangian action over a topologically constrained loop space.
This variational proof transformed the problem: rather than searching phase
space by continuation or shooting, one seeks \emph{minima of an action
functional subject to symmetry constraints}.

Sim\'o \cite{Simo2002} immediately exploited this framework numerically,
discovering hundreds of planar ($\R^2$) choreographies for small $N$ and
organising them by their discrete symmetry groups in the plane. This
catalogue constitutes the current state of the art for the two-dimensional
case. Subsequent work has extended the analysis to figure-eight stability
\cite{Chenciner2005}, choreographies with non-equal masses, and related
classes of symmetric periodic orbits (hip-hop solutions, $p/q$-choreographies).

\subsection{The three-dimensional gap}

The extension to genuine three-dimensional choreographies — orbits that
are not merely planar solutions embedded in $\R^3$ — is conspicuously
underdeveloped. The obstacles are partly technical (the symmetry group is
richer and harder to parameterise) and partly conceptual (it is less
obvious which topological sectors of the loop space admit minimisers).
The companion Mathematical Framework document \cite{MathFrame2026}
resolves the conceptual question: Propositions~4.1--4.3 and the summary
table in that document determine \emph{precisely} which three-dimensional
point groups admit genuinely non-planar choreographies and which force
the orbit into a plane. Non-planarity is not an empirical question to
be settled numerically; it is a structural consequence of the
representation theory of $G$ on the Fourier coefficient space.
We identify the systematic search over the non-planar symmetry classes
as the central programme motivating this work.

\subsection{Chemistry as organising principle}

The symmetry group of a three-dimensional choreographic orbit is a
subgroup of $O(3)$ acting on the curve $\gamma$. These are precisely the
\emph{point groups} familiar from molecular symmetry: $C_n$, $C_{nv}$,
$C_{nh}$, $D_n$, $D_{nh}$, $D_{nd}$, $T$, $T_d$, $T_h$, $O$, $O_h$,
$I$, $I_h$. The Sch\"onflies classification, originally developed to
organise molecular geometry, provides a ready-made taxonomy for three-di\-men\-sional
choreographies. Orbits with high symmetry (icosahedral $I_h$, octahedral
$O_h$) impose strong constraints on the curve $\gamma$ that dramatically
reduce the search space and may guarantee the existence of minimisers
via symmetric criticality arguments \cite{Palais1979}.

A second bridge to chemistry is topological. The spacetime trace of a
choreographic orbit — the curve $t \mapsto (\q_1(t), t)$ projected back
to $\R^3$ after closing the time axis — is a closed curve whose knot type
is a topological invariant of the orbit. Molecular knot synthesis
(trefoil knots, Solomon links, torus knots) has undergone rapid
development \cite{Stoddart2020}; a gravitational analogue in which orbit
topology is controlled by point-group symmetry would constitute a
meaningful conceptual contribution to both fields.

% ---------- 2. MATHEMATICAL FRAMEWORK ----------
\section{Mathematical Framework}

\subsection{Action functional and choreographic constraint}

The Lagrangian action over one period is
\[
  \Action[\q_1, \ldots, \q_N]
  \;=\;
  \int_0^T \Lag\,dt
  \;=\;
  \int_0^T
  \left[
    \sum_{i=1}^N \tfrac{m}{2}\,|\dot{\q}_i|^2
    \;+\;
    \sum_{i < j} \frac{Gm^2}{|\q_i - \q_j|}
  \right] dt.
\]
Restricting to the choreographic subspace via $\q_i(t) = \q_1(t +
(i-1)T/N)$ reduces the problem to a single-curve functional
\[
  \Action[\gamma]
  \;=\;
  \frac{N}{T}\int_0^T
  \left[
    \tfrac{m}{2}|\dot\gamma|^2
    + \frac{1}{N}\sum_{k=1}^{N-1}
      \frac{Gm^2}{|\gamma(t) - \gamma(t + kT/N)|}
  \right] dt,
\]
where $\gamma : [0,T] \to \R^3$ is $T$-periodic. Critical points of
$\Action$ on the appropriate Sobolev space $H^1(S^1, \R^3)$ correspond
to choreographic solutions. Global minimisers in a given homotopy class
exist whenever collisions can be excluded from that class
\cite{CM2000, Terracini2004}.

\subsection{Symmetry constraints and reduction}

Let $G \subset O(3) \times \Z/N\Z$ be the symmetry group of the orbit,
acting on the loop space by simultaneously rotating the curve and permuting
the bodies. By the Principle of Symmetric Criticality \cite{Palais1979},
a symmetric critical point of $\Action|_{\mathrm{Fix}(G)}$ is a critical
point of the full action. This reduces the search to the fixed-point
subspace $\mathrm{Fix}(G)$, whose dimension is determined by the
representation theory of $G$ on $H^1(S^1, \R^3)$.

Concretely, for a $C_{3v}$ orbit with $N=3$ bodies, the constraint forces
$\gamma$ to pass through specific planes and axes, yielding a finite-di\-men\-sional
initial-condition manifold to search. For icosahedral symmetry with
$N = 12$ or $N = 60$, the constraints are far stronger and the search space
dramatically smaller.

\subsection{Knot theory of the spacetime trace}

Define the \emph{spacetime trace} of the choreography as the closed curve
\[
  \Gamma : t \mapsto \bigl(\q_1(t),\; t \bmod T\bigr) \;\in\; \R^3 \times S^1.
\]
The isotopy class of $\Gamma$ in the complement of the other bodies'
spacetime traces is a topological invariant. For planar choreographies
this reduces to a braid word in $B_N$; for three-dimensional ones it is
a richer object, potentially a knot or link in $S^3$. We will compute
knot invariants — at minimum the Alexander polynomial and the HOMFLY-PT
polynomial — for each discovered orbit.

\subsection{Stability}

Linear stability is assessed via Floquet theory: we compute the monodromy
matrix $M = \Phi(T)$ of the linearised variational equations along the
orbit and classify eigenvalues as elliptic (on the unit circle), hyperbolic
(real), or complex-hyperbolic. An orbit is \emph{linearly stable} if all
eigenvalues lie on the unit circle. We will also compute the first few
Lyapunov exponents numerically to characterise the nonlinear stability
basin.

% ---------- 3. RESEARCH QUESTIONS ----------
\section{Research Questions}

\begin{question}
For each three-dimensional point group $G$ (Sch\"onflies classification),
does there exist a choreographic $N$-body orbit with symmetry group
containing $G$? For which $N$ and $G$ pairs do minimisers exist in the
relevant homotopy class?
\end{question}

\begin{question}
What is the knot type of the spacetime trace of each discovered orbit?
Is there a systematic map from symmetry class to knot invariant?
\end{question}

\begin{question}
What is the stability character (elliptic, hyperbolic, complex-hyperbolic)
of each discovered orbit? Are high-symmetry orbits generically more or
less stable than low-symmetry ones?
\end{question}

\begin{question}
Do the action values $\Action^\star(G, N)$ at minimisers admit a
closed-form or asymptotic expression in the symmetry order $|G|$ and
body count $N$?
\end{question}

\begin{question}
Is there a formal correspondence between orbits whose spacetime traces
realise specific knot types and known or hypothetical molecular knot
topologies?
\end{question}

% ---------- 4. METHODOLOGY ----------
\section{Methodology}

\subsection{Phase 1: Symmetry-constrained search (N = 3, 4)}

For each point group $G$ compatible with small $N$, we:
\begin{enumerate}[leftmargin=1.8em]
  \item Parameterise the fixed-point subspace $\mathrm{Fix}(G)$ using
        a Fourier truncation $\gamma(t) \approx \sum_{k=0}^{K} (a_k \cos(2\pi kt/T)
        + b_k \sin(2\pi kt/T))$ with symmetry-imposed linear constraints
        on the coefficients $a_k, b_k \in \R^3$.
  \item Minimise $\Action$ over the truncated space using gradient
        descent with an adaptive step (L-BFGS), starting from multiple
        random initialisations within the symmetry-compatible manifold.
  \item Refine converged candidates using a symplectic integrator
        (Yoshida 6th-order or Forest--Ruth) and Newton--Raphson
        correction to a true periodic orbit.
  \item Discard near-collisions and planar solutions (those for which
        $\gamma$ lies within $\varepsilon$ of a fixed plane).
\end{enumerate}

\subsection{Phase 2: Higher symmetry and larger N}

Extend to $N = 6, 12$ with tetrahedral and octahedral symmetry groups,
exploiting the stronger constraints to keep the search space tractable.
Extend further to icosahedral symmetry at $N = 12$ and $N = 60$.
The \emph{dimensional miracle} established in MathFrame \cite{MathFrame2026}
§4.6 shows that the search-space dimension at fixed truncation ratio
$K/N$ is independent of the group order $|G|$: icosahedral constraints
impose no greater computational burden than tetrahedral ones.
This formally justifies the icosahedral case as a primary target
rather than a speculative extension.

\subsection{Phase 3: Knot classification}

For each discovered orbit, compute:
\begin{itemize}[leftmargin=1.8em]
  \item The writhe and crossing number of $\Gamma$ from a projection.
  \item The Alexander polynomial $\Delta(t)$ via the Seifert matrix.
  \item The HOMFLY-PT polynomial $P(v, z)$ as a stronger invariant.
  \item The hyperbolic volume (if $\Gamma$ is hyperbolic) via SnapPy.
\end{itemize}

\subsection{Phase 4: Stability analysis}

Integrate the variational equations alongside the orbit to obtain the
monodromy matrix. Compute eigenvalues and classify stability. For
marginally stable orbits, explore KAM tori numerically via surface-of-section.

\subsection{Phase 5: Visualisation}

Produce publication-quality renders: orbital traces as tubes in $\R^3$
with point-group symmetry axes shown explicitly, spacetime braids rendered
as three-dimensional knot diagrams, and high-fidelity video renders
(4K/60fps, H.264) of the full choreographic motion for public dissemination.

% ---------- 5. EXPECTED OUTPUTS ----------
\section{Expected Outputs}

\begin{enumerate}[leftmargin=1.8em]
  \item \textbf{A catalogue of 3D choreographies}, organised by
        $(N, G, \text{knot type})$, with action values, stability
        characters, and initial conditions provided to full double
        precision.
  \item \textbf{A correspondence table} mapping symmetry class to knot
        invariant, with conjectures about which cells are empty and why.
  \item \textbf{A computational toolkit} (Python, open-source) for
        reproducing all results and searching for new orbits.
  \item \textbf{A research paper} suitable for submission to a journal
        such as \textit{Celestial Mechanics and Dynamical Astronomy},
        \textit{Nonlinearity}, or \textit{Communications in Mathematical Physics}.
  \item \textbf{A visual portfolio} for public dissemination across
        scientific social media.
\end{enumerate}

% ---------- 6. DISSEMINATION ----------
\section{Dissemination Plan}

The work will be released in two stages. The catalogue and toolkit will
be posted simultaneously to arXiv (math-ph / math.DS) and GitHub under
a permissive open-source licence, with DOI minting via Zenodo to ensure
citability. The paper will be submitted to a peer-reviewed journal
concurrently with the arXiv preprint.

Public communication will begin with the arXiv release and proceed across
a coordinated set of accounts on X, Bluesky, Farcaster, and Hacker News,
accompanied by the animated visualisations. Given the visual character of
the work and the chemistry connection, we anticipate organic reach into
both the physics and chemistry communities.

% ---------- 7. OPEN QUESTIONS ----------
\section{Open Questions and Risks}

\begin{itemize}[leftmargin=1.8em]
  \item \textbf{Collision avoidance}: Minimisers in some symmetry classes
        may pass through collisions, rendering the action infinite or
        the orbit non-physical. Excluding such classes analytically or
        verifying their emptiness numerically is a prerequisite for clean
        results.
  \item \textbf{Uniqueness}: Multiple minimisers may exist within a
        single symmetry-homotopy class. The relationship between connected
        components of $\mathrm{Fix}(G)$ and distinct orbit families is
        not understood a priori.
  \item \textbf{Fourier truncation error}: Highly irregular curves
        (low-action orbits near collision) may require large $K$ for
        accurate representation. Convergence monitoring will be required.
  \item \textbf{Genuinely new orbits}: There is a possibility — small
        but non-negligible — that 3D choreographies simply do not exist
        beyond degenerate (planar or hip-hop) cases for small $N$. A
        null result, reported carefully, would itself be a contribution.
\end{itemize}

% ---------- REFERENCES ----------
\begin{thebibliography}{9}

\bibitem{CM2000}
A.~Chenciner and R.~Montgomery,
\textit{A remarkable periodic solution of the three-body problem in the case of equal masses},
Ann.\ Math.\ \textbf{152} (2000), 881--901.

\bibitem{Simo2002}
C.~Sim\'o,
\textit{New families of solutions in $N$-body problems},
in \textit{European Congress of Mathematics}, Birkh\"auser, Basel, 2001, pp.\ 101--115.

\bibitem{Chenciner2005}
A.~Chenciner et al.,
\textit{Are there perverse choreographies?},
in \textit{New Advances in Celestial Mechanics and Hamiltonian Systems},
Springer, 2004, pp.\ 63--76.

\bibitem{Palais1979}
R.~S.~Palais,
\textit{The principle of symmetric criticality},
Commun.\ Math.\ Phys.\ \textbf{69} (1979), 19--30.

\bibitem{Terracini2004}
S.~Terracini and A.~Venturelli,
\textit{Symmetric trajectories for the $2N$-body problem with equal masses},
Arch.\ Rational Mech.\ Anal.\ \textbf{184} (2007), 465--493.

\bibitem{Stoddart2020}
F.~Schaufelberger et al.,
\textit{Synthesis of molecular knots},
Chem.\ Soc.\ Rev.\ \textbf{50} (2021), 3473--3500.

\bibitem{MathFrame2026}
e5bed97b3349,
\textit{Three-Dimensional Choreographic Orbits: Mathematical Framework},
\texttt{orbraid} repository, \texttt{docs/mathframe/mathframe.pdf} (2026).

\end{thebibliography}

\end{document}
